\section{2012-2013学年线性代数I（H）期末}

\begin{center}
    任课老师：统一命卷\hspace{4em} 考试时长：120分钟
\end{center}

\begin{enumerate}
    \item （10分）求实线性方程组 $\begin{cases}x_1-x_2+2x_3 = 1 \\ x_1-2x_2-x_3=2 \\ 3x_1-x_2+5x_3=3 \\ 2x_1-2x_2-3x_3 = 4\end{cases}$ 的解集.

    \item （10分）设 $A$ 是域 $\mathbf{F}$ 上的 $m\times n$ 矩阵，$A$ 的秩 $r(A)=1$.
    \begin{enumerate}
        \item 证明存在（列向量）$X\in \mathbf{F}^m$ 和 $Y\in \mathbf{F}^n$ 使得 $A=XY^\mathrm{T}$,其中 $Y^\mathrm{T}$ 是 $Y$ 的转置.

        \item $X$ 和 $Y$ 是否唯一？
    \end{enumerate}

    \item （10分）定义线性映射 $T\colon \mathbf{R}^n \to \mathbf{R}^n$ 如下：对任意的 $(x_1,x_2,\ldots,x_n) \in \mathbf{R}^n$
    \[T(x_1,x_2,\ldots,x_n)=(x_1+x_2,x_2+x_3,\ldots,x_n+x_1).\]
    试给出 $T$ 的核 $\ker T$ 和 $T$ 的像 $\im T$ 的维数.

    \item （10分）设 $V$ 是域 $\mathbf{F}$ 上有限维线性空间，$T\colon V\to V$ 是线性映射. 证明 $V$ 的非零向量都是 $T$ 的特征向量当且仅当存在 $\lambda \in \mathbf{F}$，使 $T(v)=\lambda v$ 对于任何 $v \in V$ 成立.

    \item （10分）设 $A=\begin{pmatrix}a & b \\ c & d\end{pmatrix}$ 是可逆矩阵，其中 $b\neq 0$；$\lambda$ 是 $A$ 的特征值.
    \begin{enumerate}
        \item 证明 $\lambda \neq 0$.

        \item 证明 $(b,\lambda-a)^\mathrm{T}$ 是属于 $\lambda$ 的特征向量.

        \item 若 $A$ 有两个不同的特征值 $\lambda_1$ 和 $\lambda_2$，求可逆矩阵 $P$ 使 $P^{-1}AP$ 是对角矩阵.
    \end{enumerate}

    \item （10分）实对称矩阵 $A=\begin{pmatrix}0 & -1 & 1 \\ -1 & 0 & 1 \\ 1 & 1 & 0\end{pmatrix}$.  求正交矩阵 $Q$ 使 $Q^\mathrm{T}AQ$ 是对角矩阵.

    \item （10分）设 $V$ 是欧式空间，$T\colon V\to V$ 是线性映射，$\lambda \in \mathbf{R}$， $u$ 是 $V$ 的非零向量. 证明：$\lambda$ 是 $T$ 的特征值且 $u$ 是属于 $\lambda$ 的特征向量当且仅当对于任何 $v\in V$ 成立 $\langle T(u),v\rangle = \lambda\langle u,v\rangle$.

    \item （10分）设 $n$ 阶实对称阵 $A$ 满足方程 $A^2-6A+5E_n=0$，其中 $E_n$ 是 $n$ 阶单位矩阵.
    \begin{enumerate}
        \item 证明 $A$ 是正定的.

        \item 若 $n=2$，试给出全部有可能与 $A$ 相似（注意，不是相合！）的对角矩阵.
    \end{enumerate}

    \item （20分）判断下面命题的真伪. 若它是真命题，给出一个简单证明；若它是伪命题，举一个具体的反例将它否定.
    \begin{enumerate}
        \item 若有限维线性空间 $V$ 的线性映射 $T\colon V \to V$ 是可对角化的，则 $T$ 是同构.

        \item 若 $A,B$ 是对称矩阵，则 $AB$ 也是对称矩阵.

        \item 若 $n$ 阶方阵 $A,B$ 中的 $A$ 是可逆的，则 $AB$ 与 $BA$ 是相似的.

        \item 若 $n>1$ 阶方阵 $A$ 的特征多项式是 $f(\lambda)=\lambda^n$，则 $A$ 是零矩阵.
    \end{enumerate}
\end{enumerate}

\clearpage
